\chapter{Communication Complexity — Det Complexity}
%%% generated by the blueprint pipeline from TCSlib.CommunicationComplexity.DeterministicCC.DetComplexity
\section{Overview}

This module defines the deterministic communication complexity of a two-party function
$f : X \to Y \to \alpha$ as the infimum of the complexity (number of bits exchanged) over
all deterministic protocols that compute $f$.  It also establishes the key
characterizations of when this complexity is at most a given natural number $n$, both
in terms of general protocols and in terms of finite-message protocols.

%%%%%%%%%%%%%%%%%%%%%%%%%%%%%%%%%%%%%%%%%%%%%%%%%%%%%%%%%%%%%%%%%%%%%%%%%%%%%%%
\section{Declarations}

\begin{theorem}[Infimum of an extended-natural family bounded by a coercion]
\lean{CommunicationComplexity.Internal.enat_iInf_le_coe_iff}
\label{CommunicationComplexity.Internal.enat_iInf_le_coe_iff}
\leanok
\difficulty{4}
Let $(f_i)_{i \in \iota}$ be a family of values in the extended natural numbers
$\bbn_\infty = \bbn \cup \{\infty\}$, indexed by an arbitrary type $\iota$, and let $n
\in \bbn$. Then $\inf_{i} f_i \le n$ if and only if there exists an index $i$ with $f_i
\le n$, where $n$ is regarded as an element of $\bbn_\infty$ via the canonical
inclusion.
\end{theorem}

\begin{definition}[Deterministic communication complexity]
\lean{CommunicationComplexity.Deterministic.communicationComplexity}
\statementsource{roughgarden-cc-bootcamp}{pdf-d54d402b16e2-p002-b002}
\statementsource{raoyehudayoff-ch01-deterministic-protocols}{pdf-fa052d8c95c0-p003-b003}
\label{CommunicationComplexity.Deterministic.communicationComplexity}
\leanok
\uses{CommunicationComplexity.Deterministic.Protocol, CommunicationComplexity.Deterministic.Protocol.Computes, CommunicationComplexity.Deterministic.Protocol.complexity}
The deterministic communication complexity of $f : X \to Y \to \alpha$ is defined as
\[
  D(f) \;=\; \inf_{\substack{p \;:\; \mathrm{Protocol}\;X\;Y\;\alpha \\ p \text{ computes } f}}
              p.\mathrm{complexity} \;\in\; \mathbb{N}_\infty.
\]
The infimum is taken over all deterministic protocols $p$ that compute $f$, and the
result lives in $\mathrm{ENat}$ to accommodate the case where no finite protocol exists.
\end{definition}

\begin{theorem}[Threshold characterization of deterministic communication complexity]
\lean{CommunicationComplexity.Deterministic.communicationComplexity_le_iff}
\label{CommunicationComplexity.Deterministic.communicationComplexity_le_iff}
\leanok
\uses{CommunicationComplexity.Deterministic.Protocol, CommunicationComplexity.Deterministic.Protocol.Computes, CommunicationComplexity.Deterministic.Protocol.complexity, CommunicationComplexity.Deterministic.communicationComplexity, CommunicationComplexity.Internal.enat_iInf_le_coe_iff}
\difficulty{1}
Let $f : X \to Y \to \alpha$ be a two-argument function and let $n$ be a natural number.
Then the deterministic communication complexity $D(f)$ satisfies $D(f) \le n$ if and
only if there exists a deterministic communication protocol $p$ over inputs $X$ and $Y$
that computes $f$ and whose communication complexity is at most $n$.
\end{theorem}

\begin{theorem}[Finite-message characterization of deterministic communication complexity]
\lean{CommunicationComplexity.Deterministic.communicationComplexity_le_iff_finiteMessage}
\label{CommunicationComplexity.Deterministic.communicationComplexity_le_iff_finiteMessage}
\leanok
\uses{CommunicationComplexity.Deterministic.FiniteMessage.Protocol, CommunicationComplexity.Deterministic.FiniteMessage.Protocol.complexity, CommunicationComplexity.Deterministic.FiniteMessage.Protocol.ofProtocol_equiv, CommunicationComplexity.Deterministic.FiniteMessage.Protocol.run, CommunicationComplexity.Deterministic.FiniteMessage.Protocol.toProtocol, CommunicationComplexity.Deterministic.FiniteMessage.Protocol.toProtocol_complexity, CommunicationComplexity.Deterministic.FiniteMessage.Protocol.toProtocol_run, CommunicationComplexity.Deterministic.communicationComplexity, CommunicationComplexity.Deterministic.communicationComplexity_le_iff}
\difficulty{4}
Let $X$, $Y$, and $\alpha$ be types, let $f : X \to Y \to \alpha$ be a two-argument
function, and let $n$ be a natural number. Then the deterministic communication
complexity satisfies $D(f) \le n$ if and only if there exists a finite-message protocol
$p$ over inputs $X$, $Y$ and outputs $\alpha$ whose execution computes $f$, in the sense
that $p$ run on inputs $x$ and $y$ returns $f(x)(y)$ for all $x$ and $y$, and whose
communication complexity is at most $n$.
\end{theorem}

\begin{theorem}[Lower-bound characterization of communication complexity]
\lean{CommunicationComplexity.Deterministic.le_communicationComplexity_iff}
\label{CommunicationComplexity.Deterministic.le_communicationComplexity_iff}
\leanok
\uses{CommunicationComplexity.Deterministic.Protocol, CommunicationComplexity.Deterministic.Protocol.Computes, CommunicationComplexity.Deterministic.Protocol.complexity, CommunicationComplexity.Deterministic.communicationComplexity}
\difficulty{1}
Let $f : X \to Y \to \alpha$ be a two-argument function, let $D(f) \in \bbn_\infty$
denote its deterministic communication complexity, and let $n$ be a natural number. Then
$n \le D(f)$ if and only if every deterministic protocol $p$ that computes $f$ has
communication complexity at least $n$.
\end{theorem}
